\section*{Chapter \ref{ch:g7:identifying-substances} --- Identifying Substances}

\begin{solution}{exo:g7:identifying-substances:1}
The limewater test: the gas is bubbled through clear limewater, which
turns milky if the gas is carbon dioxide.
\end{solution}

\begin{solution}{exo:g7:identifying-substances:2}
The powder is anhydrous copper sulfate, and the liquid contains water.
\end{solution}

\begin{solution}{exo:g7:identifying-substances:3}
Oxygen.
\end{solution}

\begin{solution}{exo:g7:identifying-substances:4}
Pencil does not \omterm{def:g2:mixing-and-dissolving:dissolve}{dissolve} in the \omterm{def:g6:solutions-and-solubility:solute}{solvent}. An ink line would itself be
carried up and separated, and spoil the \omterm{def:g7:identifying-substances:chromatography}{chromatogram}.
\end{solution}

\begin{solution}{exo:g7:identifying-substances:5}
If the spot were under the \omterm{def:g6:solutions-and-solubility:solute}{solvent}, the dyes would \omterm{def:g2:mixing-and-dissolving:dissolve}{dissolve} into the
\omterm{def:g6:solutions-and-solubility:solute}{solvent} of the beaker instead of climbing up the paper.
\end{solution}

\begin{solution}{exo:g7:identifying-substances:6}
No, it is a \omterm{def:g6:pure-substances-and-mixtures:mixture}{mixture} of two dyes: probably the yellow dye and the blue
dye of the references.
\end{solution}

\begin{solution}{exo:g7:identifying-substances:7}
Collect a little of the gas in a test tube and hold a lighted splint at
its mouth: a squeaky pop shows hydrogen.
\end{solution}

\begin{solution}{exo:g7:identifying-substances:8}
\omterm{def:g4:air-a-mixture-of-gases:air}{Air} breathed out contains much more carbon dioxide than fresh \omterm{def:g4:air-a-mixture-of-gases:air}{air}.
\end{solution}

\begin{solution}{exo:g7:identifying-substances:9}
The test shows that water is present, not that it is alone: salty water,
juice or any watery \omterm{def:g6:pure-substances-and-mixtures:mixture}{mixture} would turn the powder blue too.
\end{solution}

\begin{solution}{exo:g7:identifying-substances:10}
Three dyes (three spots). The yellow dye climbed highest.
\end{solution}

\begin{solution}{exo:g7:identifying-substances:11}
For instance: a glowing splint into each jar --- the one that relights
it holds oxygen. On the other two, a lighted splint at the mouth --- a
squeaky pop shows hydrogen. The last jar holds carbon dioxide (it can be
confirmed with limewater). Two or three tests are enough.
\end{solution}

\begin{solution}{exo:g7:identifying-substances:12}
The gas may contain no carbon dioxide, or too little to show in the
time. A blank with \omterm{def:g4:air-a-mixture-of-gases:air}{air} breathed out, which is known to contain carbon
dioxide, checks that the limewater is working.
\end{solution}

\begin{solution}{pb:g7:identifying-substances:1}
\textbf{1.} So that the conditions (paper, \omterm{def:g6:solutions-and-solubility:solute}{solvent}, time) are the same
for all four: only then can the heights of the spots be compared.

\textbf{2.} A \omterm{def:g6:pure-substances-and-mixtures:mixture}{mixture}: its ink gives three separate spots, three dyes.

\textbf{3.} Two dyes. No: the note's ink has three dyes, and pen 1 has
neither the red nor the yellow dye.

\textbf{4.} Its three spots are not at the same heights as those of the
note: they are three different dyes.

\textbf{5.} Pen 2: its three spots are exactly at the heights of the
three spots of the note.

\textbf{6.} Oxygen.

\textbf{7.} Hydrogen.

\textbf{8.} Carbon dioxide.

\textbf{9.} A blank (a control). It shows that the limewater does not
turn milky by itself, nor with ordinary \omterm{def:g4:air-a-mixture-of-gases:air}{air} in that time: so the
milkiness with gas Z really comes from carbon dioxide.

\textbf{10.} ``Water --- anhydrous copper sulfate --- white turns blue'':
the liquid contains water.

\textbf{11.} No. A positive test shows that water is present, not that
nothing else is: the liquid could be a solution or another watery
\omterm{def:g6:pure-substances-and-mixtures:mixture}{mixture}.

\textbf{12.} The ink of the note contains 3 dyes, and pen 2 wrote the
note.
\end{solution}
